
\section{Theoretical Analysis}

\subsection{Completeness}

\begin{theorem}
TA-A* is complete.
\end{theorem}

\begin{proof}
We prove completeness by showing that TA-A* systematically explores the search space:

\begin{enumerate}
\item TA-A* is based on A*, which is known to be complete. A* systematically explores the search space
\item TA-A* adds terrain evaluation to A* cost function. g(n) = g_base(n) + g_terrain(n), where g_terrain(n) >= 0
\item Non-negative terrain cost preserves admissibility. If h(n) is admissible, h_ta(n) = h(n) + terrain_cost is also admissible when terrain_cost >= 0
\item TA-A* with admissible heuristic is complete. Systematic exploration + admissible heuristic → completeness

\end{enumerate}

Therefore, if a solution exists, TA-A* will find it. \qed
\end{proof}

\subsection{Optimality}

\begin{theorem}
TA-A* with admissible heuristic finds optimal solution.
\end{theorem}

\begin{proof}
Let $f(n) = g(n) + h(n)$ where:
\begin{itemize}
\item $g(n)$ is the actual cost from start to $n$
\item $h(n)$ is an admissible heuristic
\end{itemize}

1. Define total cost: f(n) = g(n) + h(n)

2. g(n) = path_cost(start, n) + terrain_cost(start, n)

3. If h(n) is admissible: h(n) <= h*(n)

4. TA-A* never expands nodes with f(n) > C*

5. First solution found is optimal


Thus, TA-A* finds the optimal solution. \qed
\end{proof}

\subsection{Complexity Analysis}

\begin{theorem}
The time complexity of TA-A* is $O(b^d)$ where $b$ is the branching factor and $d$ is the solution depth.
\end{theorem}

Time complexity: O(b^d)

Space complexity: O(b^d)

\subsection{Convergence}

\begin{theorem}
TA-A* converges in finite steps.
\end{theorem}

\begin{proof}
1. Search space is finite

2. Each node is expanded at most once

3. Open set is non-increasing after each expansion

4. Algorithm terminates when goal is found or open set is empty


\qed
\end{proof}
